\section{Introduction}
\label{sec:intro}
When users of a platform influence one another, the effect an A/B test is meant to inform is the \emph{global} average treatment effect (GATE):
the mean outcome if everybody were treated minus the mean outcome if nobody were. The usual difference in means estimates the effect of treating a unit
\emph{while its neighbours stay untreated}, and is biased for the GATE by the spillover. A popular remedy is the regression-based estimator of
\citet{chin2019} and \citet{hanugander2023}: regress outcomes on features of the assignment vector (own treatment, the treated fraction of one's neighbours, \dots),
predict the outcomes under all-treated and all-control assignment, and difference. Its guarantees are known for a correctly specified linear model
and in the large-sample limit; how accurate it is \emph{in practice}, what the error is made of, and when it fails is less clear.

We study these questions in a controlled simulation where the truth is known exactly: we define the potential-outcome functions, so the GATE of every simulated world is computed,
not estimated, and bias and variance can be separated and attributed. We vary interference strength and shape, graph family (synthetic and two
public Facebook graphs), size and degree, assignment design, noise, and the regression specification, always against a correctly specified control, and report
Monte-Carlo uncertainty on every number. Our contributions are:
\begin{itemize}\setlength\itemsep{0pt}
\item \textbf{A bias formula for the studied specification, and its test.} For the neighbour-fraction OLS$(z,f)$, one-hop spillover, Bernoulli($p$) assignment and the noiseless population limit, the bias is $\gamma(s-1)$ with $s$ a covariance ratio over the actual degree mixture (Eq.~\ref{eq:bias}). The principle behind it (a misspecified linear fit carries the local slope to the corners) is not new and is implicit in \citet{chin2019} and \citet{hanugander2023}; what we add is the closed form for finite and heterogeneous degrees, which predicts when the sign changes with $p$, and its test: it matches 144 simulated noiseless cells (all within 3 Monte-Carlo SE), does not shrink with $n$, and is proportional to $\gamma$ in a later grid (14 of 16 groups). A leverage heuristic $\mathrm{SD}\propto\sqrt{d/(n\,p(1-p))}$ is confirmed in its exponents ($-0.501\pm0.005$ in $n$, $0.495\pm0.006$ in $d$) on Erd\H{o}s--R\'enyi graphs.
\item \textbf{A map of when the estimator is accurate}, at the grids' sizes, degrees and noise levels. It beats the difference in means as soon as interference is moderate, but pays a 5--19$\times$ larger SD when there is none;
under nonlinear spillover its error is bias-dominated and sign-reversing in $p$; richer specifications (also penalised) and boosting mostly do not help; cluster
randomisation helps only on some graphs; and intervals that ignore dependence undercover, as does a cluster bootstrap.
\item \textbf{Negative and mixed results, reported as such}: the large-degree shortcut for the bias is a poor approximation on real graphs; the predicted penalty for heavy-tailed degrees was \emph{not} observed, also in a matched test; good graph clusters \emph{hurt} the regression estimator on the dense Stanford3 graph, which replicates on a second seed although our proposed mechanism does not hold there; Hájek exposure baselines are unestimable under Bernoulli assignment at the degrees of our main grids, and where estimable (low degree, cluster designs) the pre-stated rules for them gave ``unclear'' verdicts.
\end{itemize}
All results are generated by released code from fixed seeds (Appendix~\ref{app:repro}); conclusions are scoped to the simulated regimes: model-defined truth, one outcome family per mechanism, no real platform data.
